\begin{lemma}[Equivariant extension classes for finite locally free modules]
\label{lemma-equivariant-extension-classes-finite-locally-free}
\begin{reference}
\cite[Chapter V, \S 5.6, pp. 218--219]{Tohoku}
\end{reference}
Let $(X,\mathcal O)$ be a commutative $G$-ringed space and let
$\mathcal E,\mathcal F$ be $G$-$\mathcal O$-modules. Assume that the
underlying $\mathcal O$-module $\mathcal E$ is locally isomorphic to
$\mathcal O^r$ for some $r\geq 0$. Then equivalence classes of short exact
sequences
$$
0 \longrightarrow \mathcal F \longrightarrow \mathcal V
\longrightarrow \mathcal E \longrightarrow 0
$$
in $\textit{Mod}_G(\mathcal O)$ are canonically in bijection with
$\Ext^1_{\mathcal O,G}(X;\mathcal E,\mathcal F)$. Moreover, there are
canonical isomorphisms
$$
\Ext^1_{\mathcal O,G}(X;\mathcal E,\mathcal F)
\longrightarrow
H^1\left(X;G,\SheafHom_{\mathcal O}(\mathcal E,\mathcal F)\right)
$$
and
$$
\mathcal E^\vee\otimes_{\mathcal O}\mathcal F
\longrightarrow
\SheafHom_{\mathcal O}(\mathcal E,\mathcal F),
\qquad
\mathcal E^\vee=
\SheafHom_{\mathcal O}(\mathcal E,\mathcal O),
$$
where the second arrow is the evaluation isomorphism. Consequently,
$$
\Ext^1_{\mathcal O,G}(X;\mathcal E,\mathcal F)
\cong
H^1\left(X;G,\mathcal E^\vee\otimes_{\mathcal O}\mathcal F\right).
$$
\end{lemma}

\begin{proof}
The category $\textit{Mod}_G(\mathcal O)$ is abelian and has enough
injectives. The degree-one Yoneda description of Ext in an abelian
category identifies its degree-one Ext group with equivalence classes of
degree-one Yoneda extensions, which are exactly the displayed short exact sequences. The
definition of $\Ext^1_{\mathcal O,G}$ and Lemma
\ref{lemma-equivariant-ext-properties} identify that group with
$\Ext^1_{\textit{Mod}_G(\mathcal O)}(\mathcal E,\mathcal F)$.

\medskip\noindent
Since $\mathcal E$ is locally isomorphic to $\mathcal O^r$, Lemma
\ref{lemma-equivariant-ext-locally-free}, with $n=1$, gives the first
displayed isomorphism. The internal-hom action of Lemma
\ref{lemma-equivariant-internal-hom} makes
$\mathcal E^\vee$ a $G$-$\mathcal O$-module. The evaluation map
$$
\mathcal E^\vee\otimes_{\mathcal O}\mathcal F
\longrightarrow
\SheafHom_{\mathcal O}(\mathcal E,\mathcal F),
\qquad
\lambda\otimes f\longmapsto(e\longmapsto\lambda(e)f),
$$
is $G$-equivariant by the conjugation formula in that lemma. Locally on
$X$ it is the standard isomorphism
$(\mathcal O^r)^\vee\otimes\mathcal F\cong
\SheafHom_{\mathcal O}(\mathcal O^r,\mathcal F)$, hence it is an
isomorphism globally. Applying $H^1(X;G,-)$ and composing with the first
isomorphism proves the final assertion.
\end{proof}

\begin{remark}[Analytic boundary]
This lemma records the abstract equivariant Ext bridge. Identifying
holomorphic bundles with analytic $G$-$\mathcal O$-modules, or invoking
GAGA, requires separate analytic hypotheses and is not claimed here.
\label{remark-equivariant-extension-classes-analytic-boundary}
\end{remark}
